% ================================================================
%  conteudo/revisao.tex
% ================================================================

\section{Revisão da Literatura}

Devido a orientação da disciplina, esta seção apresenta apenas os tópicos que deveriam ser abordados na revisão da literatura, fundamentando o experimento realizado.

\subsection{Leveduras e fungos filamentosos}

\begin{itemize}
    \item Conceituação de leveduras e fungos filamentosos.
    \item Principais diferenças morfológicas e fisiológicas entre esses grupos.
    \item Importância biotecnológica e industrial desses microrganismos.
    \item Exemplos de gêneros ou espécies de interesse em processos fermentativos.
\end{itemize}

\subsection{Microrganismos associados a frutas}

\begin{itemize}
    \item Microbiota nativa presente na superfície e na polpa de frutas.
    \item Fatores que influenciam a composição microbiana, como pH, maturação e condições de cultivo.
    \item Relação entre frutas e ocorrência de leveduras de interesse tecnológico.
    \item Potencial das uvas como fonte de isolamento de microrganismos.
\end{itemize}

\subsection{Meios de cultura sólidos}

\begin{itemize}
    \item Função dos meios de cultura sólidos no isolamento microbiano.
    \item Importância do ágar na individualização e observação de colônias.
    \item Características de meios utilizados para cultivo de fungos e leveduras.
    \item Critérios para escolha do meio de cultura conforme o microrganismo-alvo.
\end{itemize}

\subsection{Técnicas de isolamento microbiano}

\begin{itemize}
    \item Princípios do isolamento de culturas puras a partir de amostras mistas.
    \item Técnica de esgotamento por estrias em placas de Petri.
    \item Coleta por swab como método de recuperação da microbiota superficial.
    \item Importância das condições assépticas durante a inoculação e incubação.
    \item Influência do tempo e da temperatura de incubação no desenvolvimento das colônias.
\end{itemize}
